\chapter{Blood Tests A--Z}

\section*{Definition and Overview}
Blood tests are laboratory investigations performed on a sample of blood. They are organised into many individual tests, each measuring a different component of the blood or a different chemical within it. This chapter provides an alphabetical guide to the more common blood tests a person may encounter -- from antibody tests to white cell counts -- explaining in plain language what each one measures and why it might be ordered. The list is not exhaustive, but it covers the tests most frequently requested in primary care and hospital practice.

\section*{Epidemiology}
Different blood tests are common at different ages and in different conditions. Iron studies are frequently requested for women of childbearing age, HbA1c for people with or at risk of diabetes, and prostate-specific antigen (PSA) for older men. Full blood counts and electrolyte panels are among the most commonly ordered tests across all settings, and thyroid and vitamin B12 tests are performed very frequently in general practice. The appropriate use of these tests supports the diagnosis and monitoring of some of the most prevalent health conditions, including anaemia, diabetes, thyroid disease, and chronic kidney disease.

\section*{Aetiology and Pathophysiology}
Selected tests and what they measure:

\begin{itemize}
    \item \textbf{A:} antibodies, such as those for hepatitis or coeliac disease; also alpha-fetoprotein and autoimmune markers
    \item \textbf{B:} bilirubin (liver function and red-cell breakdown); vitamin B12 and folate (anaemia and nerve health)
    \item \textbf{C:} calcium, cholesterol, C-reactive protein (CRP), and cortisol (bone, heart, inflammation, and stress responses)
    \item \textbf{D:} vitamin D; and a blood count differential, which separates the white cell subtypes
    \item \textbf{F:} ferritin and iron studies -- body iron stores and iron deficiency
    \item \textbf{F:} full blood count -- red cells, white cells, and platelets
    \item \textbf{G:} glucose, and HbA1c, which reflects average blood sugar over the previous months
    \item \textbf{K:} kidney function (creatinine and urea) and potassium
    \item \textbf{L:} liver function tests and lipids (cholesterol and triglycerides)
    \item \textbf{M:} magnesium and other minerals; also blood cultures that look for infection in the blood
    \item \textbf{P:} platelets and a pregnancy test
    \item \textbf{R:} renal function and a reticulocyte count
    \item \textbf{S:} sodium, and the sedimentation rate (ESR), an inflammation marker
    \item \textbf{T:} thyroid function (TSH and free T4) and troponin, a marker of heart muscle damage
    \item \textbf{U:} uric acid, associated with gout
    \item \textbf{W:} white cell count, raised in infection and some blood conditions
\end{itemize}

\section*{Clinical Features}
Most blood tests feel the same at the point of collection, whatever is being measured:

\begin{itemize}
    \item A brief sting as the needle is inserted
    \item Minor bruising or tenderness at the site
    \item Light-headedness in some people
    \item Certain tests need special preparation, such as fasting or timed blood collection
    \item Some results, such as troponin or blood cultures, are time-sensitive and collected urgently
\end{itemize}

\section*{Differential Diagnosis}
The same test can be ordered for very different reasons, and the same result can have different meanings:

\begin{itemize}
    \item A raised white cell count may reflect infection, inflammation, stress, or a blood disorder
    \item Low iron can be due to diet, heavy menstrual loss, or hidden bleeding in the gut
    \item Abnormal thyroid tests can point to overactivity or underactivity, or may be temporary
    \item An abnormal result may reflect a medication a person is taking rather than disease
    \item Results are interpreted in the context of symptoms, examination, and other tests
\end{itemize}

\section*{When to Seek Medical Care}
Contact a doctor if:

\begin{itemize}
    \item You have been told a test result is abnormal but do not know what it means
    \item You have symptoms such as unexplained tiredness, weight change, or bleeding, and tests have been suggested
    \item You are due to repeat a test and have not yet had the opportunity
    \item You experience ongoing discomfort, swelling, or infection at a test site
\end{itemize}

\textbf{Call 000 (or local emergency services) for:}
\begin{itemize}
    \item Collapse or prolonged fainting after a blood test
    \item Chest pain, difficulty breathing, or any sudden, severe symptom
\end{itemize}

\section*{Diagnosis and Investigations}
\begin{itemize}
    \item \textbf{Clinician review:} the doctor explains which tests are being ordered and why
    \item \textbf{Combined profiles:} tests such as liver function or thyroid panels are run together to give a fuller picture
    \item \textbf{Reference ranges:} results are compared with the laboratory's normal range, which varies with age, sex, and pregnancy
    \item \textbf{Serial testing:} repeating a test over time often reveals trends that are more meaningful than a single value
    \item \textbf{Confirmatory tests:} an abnormal screening test is usually followed by a more specific confirmatory test
    \item \textbf{Specialist interpretation:} haematologists, endocrinologists, and other specialists help interpret complex results
\end{itemize}

\section*{Management}
\begin{itemize}
    \item Results guide the next step: lifestyle change, medication, further investigation, or referral
    \item Some abnormal results need immediate action, while others are simply monitored
    \item Results are explained to the patient, and questions are encouraged
    \item Follow-up testing is arranged at appropriate intervals for chronic conditions
    \item If results do not match the clinical picture, the test may be repeated or reconsidered
\end{itemize}

\section*{Complications}
\begin{itemize}
    \item Bruising, bleeding, or infection at the collection site
    \item Fainting during or shortly after the test
    \item Anxiety provoked by abnormal or borderline results
    \item False-positive results leading to unnecessary worry and further testing
    \item False-negative results providing false reassurance
    \item The small harms of cumulative testing, including blood loss in the very unwell
\end{itemize}

\section*{Prognosis and Outcomes}
For most people, any single blood test is a safe, brief procedure with a quick recovery. The outcome that matters is what the result leads to -- most abnormal results are followed up with further assessment that clarifies the situation. Many common conditions detected by blood tests, including anaemia, diabetes, and thyroid disease, respond well to treatment when found in time. Regular monitoring through blood tests helps keep chronic conditions stable and catches complications earlier.

\section*{Prevention and Self-Care}
\begin{itemize}
    \item Follow preparation instructions such as fasting where these are given
    \item Understand which tests have been ordered and why, by asking questions
    \item Keep a record of results and the dates they were taken
    \item Remind your doctor of your medicines and supplements before testing
    \item Maintain routine health checks so that important conditions are not missed
    \item Report any new symptoms rather than waiting for the next scheduled test
\end{itemize}

\section*{Sources and Further Reading}
The following sources provide reliable, up-to-date information on blood tests and their interpretation:

\begin{itemize}
    \item NHS. \textit{NHS guidance on blood tests and their meaning.} https://www.nhs.uk/conditions/
    \item Mayo Clinic. \textit{Information on blood tests and results.} https://www.mayoclinic.org/
    \item MedlinePlus. \textit{Blood tests patient information.} https://medlineplus.gov/
    \item MSD Manual. \textit{Overview of laboratory tests.} https://www.msdmanuals.com/
    \item Centers for Disease Control and Prevention. \textit{Laboratory testing information.} https://www.cdc.gov/
    \item NICE. \textit{Diagnostic test guidance.} https://www.nice.org.uk/
\end{itemize}
